\section{The ray-family explicit formula and its exceptional zero}
\label{sec:ray-explicit}

Use the family constructed in Section~\ref{sec:quadratic-ray-family}.
For a character \(\chi\) write
\[
 b_\chi(n)=\sum_g\chi(g)\lambda_g(n),\qquad
 {\cal L}_\chi(z)=\sum_{n\geq1}b_\chi(n)n^{-z},\qquad
                \ell(t)=\log(D(|t|+2)).
\]
Equation~\eqref{eq:ray-lambda-budget} gives
\(|b_\chi(n)|\leq2\Lambda(n)\).
The ideal Euler product, absolutely convergent for \(\Re z>1\), gives
\[
 {\cal L}_\chi(z)=-\frac{L'_\chi(z)}{L_\chi(z)}
       ={\bf1}_{\chi=1}\frac1{z-1}-\frac{F'_\chi(z)}{F_\chi(z)}.
\]
To check convergence, the counting estimates give a total of
\(O(Du)\) ideals of norm at most \(u\).
The logarithm of the product is also dominated by
\(2\sum_{p,j}p^{-j\Re z}/j\).
Its exponential is nonzero, so every zero of \(F_\chi\) in
\(\Re z>1/2\) has real part at most one.

\subsection{A local formula with uniform constants}

There is an absolute \(C>0\) such that the zeros of \(F_\chi\)
in the disk
\[
             |z-(5/4+it)|\leq3/5
\]
have total multiplicity at most \(C\ell(t)\), and, for
\(|z-(5/4+it)|\leq1/2\) away from the zeros,
\begin{equation}
 \label{eq:ray-local-log}
 \left|\frac{F'_\chi(z)}{F_\chi(z)}
             -\sum_\rho\frac{m_\rho}{z-\rho}\right|
                                  \leq C\ell(t).
\end{equation}
We prove the assertion for the regularized principal function too.
On \(\Re z\geq11/20\), the summatory integral used above gives
\[
                    |F_\chi(z)|\leq300D(1+|z|)^2.
\]
For nontrivial characters use
\(|I_\chi(z)|\leq10D/(\Re z-1/2)\).
For the trivial character use
\(F_1(z)=\rho z+(z-1)zI_1(z)\), \(\rho\leq4D\).
At the center \(z_0=5/4+it\), the Euler product gives
\[
 |L_\chi(z_0)|\geq\zeta(5/4)^{-2},\qquad
                  |F_\chi(z_0)|\geq\frac14\zeta(5/4)^{-2}.
\]
Indeed the absolute value of the logarithm of the Euler product is
at most \(2\log\zeta(5/4)\); the extra factor in \(F_1\) has norm
at least \(1/4\).
The ratio of the upper bound on the disk of radius \(7/10\) to
this lower bound is \(M_0\) with
\(\log M_0\leq C_1\ell(t)\), for an absolute \(C_1\).

Here is the disk argument. Jensen's formula bounds the number \(k\)
of zeros in the disk of radius \(3/5\), with multiplicity, by
\(\log M_0/\log(7/6)\).
Boundary zeros can be handled by taking radii increasing to \(7/10\).
Let \(P\) be the product of the corresponding monic linear factors,
and set \(g=F_\chi/P\), extending across those zeros.
On the outer circle, \(|P|\geq(1/10)^k\), while
\(|P(z_0)|\leq(3/5)^k\). The maximum modulus principle gives
\[
                  |g(z)/g(z_0)|\leq M_0\,6^k .
\]
The function \(g\) has no zeros in the open disk of radius \(3/5\).
Its normalized logarithm \(h\) there has \(h(z_0)=0\) and
\(\Re h\leq A:=\log M_0+k\log6\).
For \(a>A\), Schwarz's lemma applied to \(h/(2a-h)\) gives
\[
                  |h(z)|\leq\frac{2a|z-z_0|}{3/5-|z-z_0|}.
\]
On the disk of radius \(11/20\) this is at most \(22a\).
Cauchy's derivative estimate on a circle of radius \(1/20\)
about any point in the disk of radius \(1/2\), followed by
\(a\downarrow A\), gives \(|h'|\leq440A\).
This proves \eqref{eq:ray-local-log} and the asserted zero count.

Since \(\Re\rho\leq1\), every zero term in
\eqref{eq:ray-local-log} has nonnegative real part at
\(s=1+a+it\), \(a>0\). Dropping unused zero terms gives, for
\(0<a\leq1/2\),
\begin{equation}
 \label{eq:ray-real-log-upper}
 \Re{\cal L}_\chi(1+a+it)
 \leq{\bf1}_{\chi=1}\frac{a}{a^2+t^2}
       +C\ell(t)-\sum_{\rho\in S}
                m_\rho\Re\frac1{1+a+it-\rho},
\end{equation}
for any selection \(S\) of the local zeros. Enlarge \(C\), when
necessary, so that the same constant works throughout.

\subsection{The zero-free strip}

Positivity of the prime weights gives
\begin{equation}
 \label{eq:ray-341}
 3{\cal L}_1(\sigma)+4\Re{\cal L}_\chi(\sigma+it)
                        +\Re{\cal L}_{\chi^2}(\sigma+2it)\geq0
                        \qquad(\sigma>1).
\end{equation}
Termwise this is \(3+4\cos u+\cos2u=2(1+\cos u)^2\geq0\).
For real \(\chi,\psi\) it also gives
\[
 {\cal L}_1(\sigma)\geq0,\qquad
 {\cal L}_1(\sigma)+{\cal L}_\chi(\sigma)\geq0,
\]
and
\[
 {\cal L}_1(\sigma)+{\cal L}_\chi(\sigma)
       +{\cal L}_\psi(\sigma)+{\cal L}_{\chi\psi}(\sigma)\geq0.
\]
The last two assertions follow respectively from \(1+\chi(g)\geq0\)
and \((1+\chi(g))(1+\psi(g))\geq0\).

Suppose \(\rho=1-\delta+i\gamma\) is a zero,
\(0\leq\delta\leq\epsilon\leq1/400\), with
\(\chi^2\ne1\) or \(\gamma\ne0\).
We show
\begin{equation}
 \label{eq:ray-epsilon}
                           \epsilon^{-1}\leq300C\ell(\gamma).
\end{equation}
First suppose that \(\chi^2\ne1\) or \(|\gamma|\geq30\epsilon\).
Use \(a=6\epsilon\) in
\eqref{eq:ray-real-log-upper} and \eqref{eq:ray-341}.
The selected zero contributes at least \(1/(7\epsilon)\).
Every possible principal pole at \(\gamma\) or \(2\gamma\) is at
most \(1/(156\epsilon)\), since
\[
          \frac{6\epsilon}{(6\epsilon)^2+\gamma^2}
                         \leq\frac1{156\epsilon}
                         \quad(|\gamma|\geq30\epsilon).
\]
When \(\chi^2\ne1\), both of those possible pole terms vanish.
Also \(\ell(0)\leq\ell(\gamma)\) and
\(\ell(2\gamma)\leq2\ell(\gamma)\).
Thus
\[
 0\leq
 \left(\frac12-\frac47+\frac5{156}\right)\epsilon^{-1}
                                      +9C\ell(\gamma)
       =-\frac{43}{1092}\epsilon^{-1}+9C\ell(\gamma),
\]
which implies \eqref{eq:ray-epsilon}.

In the other case \(\chi\) is real and
\(0<|\gamma|<30\epsilon\).
Conjugation gives the distinct second zero \(1-\delta-i\gamma\).
Both zeros lie in the local disk centered at \(5/4\).
Use \(a=60\epsilon\).
If \(\chi=1\), apply \({\cal L}_1(\sigma)\geq0\) directly.
If \(\chi\ne1\), apply
\({\cal L}_1(\sigma)+{\cal L}_\chi(\sigma)\geq0\), using the
empty zero selection for the first term.
In either case \eqref{eq:ray-real-log-upper} gives
\[
 2\frac{60\epsilon+\delta}{(60\epsilon+\delta)^2+\gamma^2}
                     \leq\frac1{60\epsilon}+2C\ell(0).
\]
The fraction on the left before multiplication by \(2\) is at
least \(60/(4621\epsilon)\). Subtracting \(1/(60\epsilon)\)
gives
\[
           \frac{2579}{277260}\epsilon^{-1}
                                      \leq2C\ell(0),
\]
which again implies \eqref{eq:ray-epsilon}.

Set \(K=300C\) and \(c_1=(4K+400)^{-1}\).
If such a zero had
\(\delta\leq c_1/\ell(\gamma)\), take
\(\epsilon=\{(2K+200)\ell(\gamma)\}^{-1}\).
Then \(0<\epsilon\leq1/400\), \(\delta\leq\epsilon\), and
\eqref{eq:ray-epsilon} would give \(2K+200\leq K\).
Consequently every zero in
\(\Re z\geq1-c_1/\ell(\Im z)\) is real and belongs to a real
character.

We also prove uniqueness and simplicity across that real family.
Choose
\[
                 c_2=\frac1{100(C+1)},\qquad
                    \sigma=1+\frac{2c_2}{\ell(0)}.
\]
A zero \(\beta\geq1-c_2/\ell(0)\), of multiplicity \(m\), contributes
at least \(m\ell(0)/(3c_2)\) to the zero sum at \(\sigma\).
If one character had two such zeros, or one multiple zero, choose
two zeros counting multiplicity.
Use \({\cal L}_1\geq0\) when that character is trivial, and
\({\cal L}_1+{\cal L}_\chi\geq0\) otherwise.
If two different real characters had such zeros, use
\({\cal L}_1+{\cal L}_\chi\geq0\) when one is trivial, and the
four-term positive expression when both are nontrivial.
In the last case \(\chi\psi\ne1\).
There is one principal pole in each expression. Its size is
\(\ell(0)/(2c_2)\), and the combined error is at most \(4C\ell(0)\).
Thus every case would imply
\[
             \frac{2}{3c_2}\leq\frac1{2c_2}+4C,
                         \qquad \frac1{6c_2}\leq4C ,
\]
contrary to the choice of \(c_2\).
There is therefore at most one pair \((\chi,\beta)\) in this
strip, and its zero is simple.

Fix
\[
                    c_0=\min(c_1,c_2,1/20)>0 .
\]
Let \(Z_\chi\) contain the real zero with
\(\beta\geq1-c_0/\ell(0)\) when it exists, and let it be empty
otherwise.
It lies in \([3/4,1)\).
Section~\ref{sec:ray-real-gap} gives
\(1-\beta\geq(c_*/18000)D^{-46}\), with an absolute \(c_*>0\).

\subsection{Removing the exceptional pole}

Uniformly for
\[
       1-\frac{c_0}{2\ell(\Im s)}\leq\Re s\leq3 ,
\]
the holomorphic continuation of
\[
     A_\chi(s)=-\frac{F'_\chi(s)}{F_\chi(s)}
                           +\sum_{\beta\in Z_\chi}\frac1{s-\beta}
\]
satisfies
\begin{equation}
 \label{eq:ray-corrected-log}
                         |A_\chi(s)|\leq C_2\ell(\Im s)^2 ,
\end{equation}
for an absolute \(C_2\), depending on the fixed \(c_0\).
The apparent singularity at the simple exceptional zero is removable.
Here are the distance estimates proving the bound.
When \(\Re s\leq7/4\), use \eqref{eq:ray-local-log}.
A local zero has \(|\Im\rho-\Im s|\leq3/5\), and
\[
       \ell(\Im\rho)\leq\frac43\ell(\Im s).
\]
Indeed
\(|\Im\rho|+2\leq(13/10)(|\Im s|+2)\), and
\(\log(13/10)\leq\ell(\Im s)/3\) since \(D\geq4\).
Every local zero other than the exceptional one either has real
part below \(3/4\), or satisfies
\[
        \Re\rho<1-\frac{3c_0}{4\ell(\Im s)}.
\]
Its distance from \(s\) is therefore at least
\(c_0/(4\ell(\Im s))\).
The local zero count is \(O(\ell(\Im s))\), so their total
logarithmic derivative contribution is \(O(\ell(\Im s)^2)\).
If the exceptional zero lies outside the local disk, its distance
from \(s\) is at least \(1/10\).
For \(\Re s\geq7/4\), absolute convergence of the Euler series
bounds \({\cal L}_\chi(s)\) by a constant, and
\(|s-1|^{-1}\), \(|s-\beta|^{-1}\leq4/3\).
This proves \eqref{eq:ray-corrected-log} everywhere stated.

\subsection{Perron inversion and the contour shift}

We recall the precise truncation estimate needed here, for arbitrary
coefficients with \(|b(n)|\leq2\Lambda(n)\).
For \(u\geq2\), \(2\leq T\leq u\), and \(c=1+1/\log u\),
\[
 \left|\sum_{n\leq u}b(n)
     -\frac1{2\pi i}\int_{c-iT}^{c+iT}
            \left(\sum_nb(n)n^{-s}\right)\frac{u^s}{s}\,ds\right|
                    \leq C_3\{u(\log u)^2/T+\log u\}.
\]
For clarity, the integer nearest the cutoff requires the last term.
The truncated step kernel
\((2\pi i)^{-1}\int_{c-iT}^{c+iT}y^s\,ds/s\) differs from
\({\bf1}_{y>1}\), for \(y\ne1\), by at most
\[
                  C y^c\min(1,(T|\log y|)^{-1}).
\]
The reciprocal bound follows by integration by parts in the
oscillatory tail. The bound by a constant follows by writing its
real integral in terms of \(\sin(t\log y)\), splitting at
\(t=|\log y|^{-1}\), and applying the Dirichlet bound thereafter.
At \(y=1\) the kernel is \(\pi^{-1}\arctan(T/c)\).
Summing the errors for \(n\leq u/2\) and \(n\geq2u\) uses
\(\sum(\log n)n^{-c}=O((c-1)^{-2})\).
In \([u/2,2u]\), remove the at most two integers nearest \(u\),
whose contribution is \(O(\log u)\).
The remaining errors are bounded by a harmonic sum times
\(C(u/T)\log u\). This gives the stated uniform estimate,
including integral and nonintegral \(u\).

Set \(L=\sqrt{\log u}\) and suppose \(\log D\leq L\).
Take \(a=c_0/12\), \(T=e^{aL}\), and
\[
                  \sigma=1-\frac{c_0}{2\log(D(T+2))}.
\]
For all sufficiently large \(u\), \(2\leq T\leq u\), and
\(\log(D(T+2))\leq3L\).
The corrected function \(A_\chi\) is holomorphic on and near the
rectangle from \(\sigma-iT\) to \(c+iT\).
Shift its Perron integral to the left boundary. Equation
\eqref{eq:ray-corrected-log} bounds the horizontal contributions
by \(C uL^2/T\), and the vertical contribution by
\[
                C u^\sigma L^2\log(T+2)
                         \leq C uL^3e^{-c_0L/6}.
\]
The integral of \(1/|\sigma+it|\) is \(O(\log(T+2))\);
also \(u^c=eu\), and \(\sigma>3/4\).
All these estimates are uniform in the character and modulus.

Finally, for \(3/4\leq b\leq1\) and \(c>b\), elementary inversion
gives
\[
 \frac1{2\pi i}\int_{c-iT}^{c+iT}
                         \frac{u^s}{s(s-b)}\,ds
               =\frac{u^b}{b}+O(1+u/T),
\]
uniformly in \(b\).
The full line integral equals \((u^b-1)/b\), by the residues at
\(b\) and zero. Its two tails are \(O(u^c/T)\), because
\(|s(s-b)|\geq t^2\).
Apply this once with \(b=1\) for the principal pole and, when
present, once with \(b=\beta\) for the exceptional pole.
Since
\({\cal L}_\chi=A_\chi+{\bf1}_{\chi=1}/(s-1)
-\sum_{\beta\in Z_\chi}1/(s-\beta)\),
Perron inversion and the rectangle estimate imply
\[
 \left|\sum_{n\leq u}b_\chi(n)-{\bf1}_{\chi=1}u
               +\sum_{\beta\in Z_\chi}\frac{u^\beta}{\beta}\right|
                         \leq C_4u e^{-\eta\sqrt{\log u}}
\]
for some absolute \(C_4,\eta>0\).
For example, after increasing \(C_4\), any fixed
\(0<\eta<a\) absorbs the powers of \(L\).
The bounded range where \(T<2\) is covered by increasing the same
constant, using \(|b_\chi(n)|\leq2\log n\).
This proves \eqref{eq:level-psi-input} with its required uniformity
and unique possible exceptional zero.

Together with the preceding two sections this establishes
\eqref{eq:level-short-norm}. Thus the analytic norm-prime input to
the simultaneous auxiliary-prime construction has now been derived.

\paragraph{Source map.}
The local disk formula is
\srcid{RayClassLFunctionLogDerivLocalFormula}.
The two zero arguments are \srcid{RayClassLFunctionZeroFreeRegion}
and \srcid{RayClassLFunctionUniqueSimpleExceptionalZero}.
The corrected derivative estimate and Perron shift are
\srcid{LogDerivSubExceptionalLe},
\srcid{TruncatedPerronOfLeTwoVonMangoldt}, and
\srcid{RayClassFamilyPsiWithExceptionalZero}.
\footnote{\href{https://github.com/Deicyde/riemann-hypothesis/blob/40e56bab026552e089d7afcfeff08d2defb9c005/src/Development/AbelianVarietyTheory/Proofs/GlobalInterpolation/RayClassLFunctionZeroFreeRegion.lean}{The zero-free argument};
\href{https://github.com/Deicyde/riemann-hypothesis/blob/40e56bab026552e089d7afcfeff08d2defb9c005/src/Development/AbelianVarietyTheory/Proofs/GlobalInterpolation/RayClassLFunctionUniqueSimpleExceptionalZero.lean}{uniqueness and simplicity};
\href{https://github.com/Deicyde/riemann-hypothesis/blob/40e56bab026552e089d7afcfeff08d2defb9c005/src/Development/AbelianVarietyTheory/Proofs/GlobalInterpolation/RayClassFamilyPsiWithExceptionalZero.lean}{the explicit formula}.}
